\chapter{Conclusion and next validated steps}

This work establishes a corrected experimental foundation for ACE, not a final
universal ranking of schedulers. The legacy per-core serialization was not a
valid congestion model: it hid simultaneous requests and made the memory budget
unable to bind. The corrected implementation gives each qubit an ordered chain
of hops while allowing unrelated qubits to request EPR pairs concurrently,
uses distinct reservation identities, admits requests event-by-event, and checks
the resulting event trace with explicit invariants. The common latency definition
now includes the waiting that a request actually experiences after it becomes
eligible, while makespan remains a separate parallel-work metric.

The completed Golden slices already provide two precise observations. On GHZ and
BV with T-HQA placement and mesh--XY routing, Dynamic full-chain reaches its
best observed latency with two compiler memories and zero observed compiler or
protocol waste. CGP and ACGP can approach or tie the latency only at much larger
compiler allocations on some capacities, and their delivered fidelity and waste
must be reported alongside latency. These are validated workload-specific
results; they do not justify saying that Dynamic wins on every workload or every
memory budget. The remaining base slices, routing arms, torus arm and MQT Bench
QPE extension must complete before a cross-workload conclusion is written.

The QPE extension is deliberately independent of the ten-workload denominator.
It uses MQT Bench's factory-generated \texttt{qpeexact} circuit, fresh HQA and
T-HQA placement contracts, the same route-policy set, the same memory grid and
the same paired physics seeds. Its execution is resumable and source-locked;
interim rows are evidence of progress, not a completed QPE comparison.

\section{Figure, table and data-file rule}
The report's existing Golden figures obey the following audit trail. Each row
below names the visual family, its corresponding complete numeric table and the
machine-readable source. The same rule will be applied to every future workload.

\begin{table}[H]
\centering\small
\begin{tabular}{@{}p{0.24\linewidth}p{0.32\linewidth}p{0.34\linewidth}@{}}
\toprule
Visual family & Corresponding table & Repository data source \\ \midrule
GHZ mesh--XY envelope & Table~\ref{tab:golden-ghz-xy-envelope} & \href{https://github.com/sanjeevkrishnaa/Compiler-Driven-ACE/blob/codex/ace-scheduler-correctness/results/golden_study_2026-10-25/summaries/ghz_thqa_mesh_xy_memory_sweep_observed_minima.csv}{GHZ CSV} \\
GHZ all allocations & Table~\ref{tab:golden-ghz-xy-envelope} & \href{https://github.com/sanjeevkrishnaa/Compiler-Driven-ACE/blob/codex/ace-scheduler-correctness/results/golden_study_2026-10-25/summaries/base-evaluation-progress.csv}{progress CSV} \\
BV mesh--XY envelope & Table~\ref{tab:golden-bv-xy-envelope} & \href{https://github.com/sanjeevkrishnaa/Compiler-Driven-ACE/blob/codex/ace-scheduler-correctness/results/golden_study_2026-10-25/summaries/bv_thqa_mesh_xy_memory_sweep_observed_minima.csv}{BV CSV} \\
BV all allocations & Table~\ref{tab:golden-bv-xy-envelope} & \href{https://github.com/sanjeevkrishnaa/Compiler-Driven-ACE/blob/codex/ace-scheduler-correctness/results/golden_study_2026-10-25/summaries/base-evaluation-progress.csv}{progress CSV} \\
\bottomrule
\end{tabular}
\caption{Figure-to-table-to-CSV audit registry. A future figure must add a row here before it is treated as a result.}
\label{tab:figure-data-registry}
\end{table}

\section{What comes next}
The next scientifically valid step is completion and validation of the frozen
matrix, followed by generated cross-workload tables and figures for request
latency, p95 latency, makespan, admission wait, readiness/hit rate, pair waste,
fidelity, retries and memory occupancy. Only after those artifacts pass the same
acceptance checks should adaptive congestion thresholds, alternative routing
heuristics, or additional QPE variants be compared. Any change to physical
parameters or scheduler semantics creates a new study identifier and a new
source lock; it must not overwrite the current evidence.
